\section{Methods}
\label{app:methods}

This appendix records the numerical methods behind the results: the Fisher assembly and its
preconditioning, the Gram-rank threshold used by the classifier, the projector-based
classification algorithm, the $\{s,\dtau\}$ Schur complement, the independence of the
numpy/LAPACK oracle, and the derivation of the relativistic frame$\to$rate Jacobian. All are
consistent with the vendored provenance notice (Section~\ref{sec:availability}).

\subsection{Fisher assembly and preconditioning}
\label{app:fisher}

Each observation block $(\{\bm g_i\},\sigma)$ contributes
$\sigma^{-2}\sum_i\bm g_i\bm g_i^\top$ to the $9\times9$ Fisher \eqref{eq:fisher}, assembled
by \code{assemble\_coupled\_info}. Before accumulation, each row is preconditioned by
dividing columns $3$--$6$ (scale and the three rotations) by $\rmoon=\SI{1737400}{m}$,
bringing all nine partials to $\mathcal{O}(1)$. This is a positive diagonal column scaling
$\Imat\mapsto\bm D\Imat\bm D$; by Lemma~\ref{lem:precond} it preserves the null-space
dimension and the classification, changing only the eigenvalue \emph{magnitudes} (and hence
the condition number). All eigenvalues, the null space, and the marginal reported in the
tables are computed on the preconditioned matrix, and by Lemma~\ref{lem:precond} equal those
of the physical information.

\subsection{The Gram-rank threshold}
\label{app:gram}

The classifier needs the ranks of the row-blocks $\Umat_{\Ssp}$ and $\Umat_{\Tsp}$ of the
orthonormal null basis (Definition~\ref{def:defect}). It computes each as the number of
eigenvalues of the Gram matrix $\bm A\bm A^\top$ exceeding an absolute threshold $10^{-9}$,
where $\bm A$ is the block. Because $\Umat$ has orthonormal \emph{columns}, the Gram
eigenvalues of any row-block lie in $[0,1]$, so $10^{-9}$ cleanly separates zero from
$\mathcal{O}(1)$ and the threshold is dimensionless and invariant under positive column
rescaling. The overall defect $d$ is
taken from \code{fim::crlb} with a \emph{relative} tolerance $\text{rel\_tol}=10^{-9}$: an
eigenvalue of the full Fisher below $\text{rel\_tol}\cdot\lambda_{\max}$ is treated as zero.
For the \code{single\_epoch} network this gives a threshold
$10^{-9}\times44.97\approx4.5\times10^{-8}$, six orders below the smallest nonzero eigenvalue
$0.1605$---a clean gap (Section~\ref{sec:networks}). A genuine near-defect with a
sub-threshold temporal projection would sit on the spatial/coupled boundary; the anchored
networks do not, so the counts are unambiguous.

\subsection{The projector-based classification algorithm}
\label{app:classify}

Given the Fisher, \code{classify\_null\_space} (i)~calls \code{fim::crlb} to obtain the
orthonormal null basis $\Umat$ (layout $\Umat[\text{row}][\text{col}]$) and the defect $d$,
returning the all-zero class if $d=0$; (ii)~forms $\Umat_{\Ssp}=\Umat[0{:}7,:]$ and
$\Umat_{\Tsp}=\Umat[7{:}9,:]$ and computes $r_{\Ssp}=\rank(\Umat_{\Ssp})$,
$r_{\Tsp}=\rank(\Umat_{\Tsp})$ by the Gram threshold of Appendix~\ref{app:gram};
(iii)~derives $\dim(\Nsp\cap\Ssp)=d-r_{\Tsp}$, $\dim(\Nsp\cap\Tsp)=d-r_{\Ssp}$, and
$\operatorname{coupled\,dim}=r_{\Ssp}+r_{\Tsp}-d$ (a debug assertion guards
$r_{\Ssp}+r_{\Tsp}\ge d$ against underflow); and (iv)~forms the off-diagonal block of the
projector $\Pnull=\Umat\Umat^\top$ restricted to rows $0$--$6$ and columns $7$--$8$,
$\Pst[i][j]=\sum_{c}\Umat[i][c]\,\Umat[j][c]$, and returns
$\lVert\Pst\rVert_{\mathrm F}$. By Proposition~\ref{prop:invariance} every output depends only
on $\Nsp$, not on the eigensolver's basis choice; the unit test
\code{classify\_is\_invariant\_under\_null\_basis\_rotation} rotates $\Umat$ by
$\theta=0.6$~rad and confirms an identical class.

\subsection{The \texorpdfstring{$\{s,\dtau\}$}{\{s,dtau\}} Schur complement}
\label{app:schur}

The coupling-1 marginal (Section~\ref{sec:coupling1}) eliminates the seven parameters
$M=\{0,1,2,4,5,6,8\}$ (translations, rotations, and the rate) to expose the $2\times2$
information on $K=\{3,7\}=\{s,\dtau\}$. With the block partition of the Fisher into
$\Imat_{KK},\Imat_{KM},\Imat_{MM}$, the Schur complement is
\begin{equation}
  \bm S \;=\; \Imat_{KK}-\Imat_{KM}\,\Imat_{MM}^{+}\,\Imat_{KM}^\top,
  \label{eq:schur}
\end{equation}
where $\Imat_{MM}^{+}$ is the Moore--Penrose pseudo-inverse, taken from the
\code{fim::crlb} pseudo-covariance field, which remains finite even when the
$M$-block is rank-deficient (sparse networks). The complement \eqref{eq:schur} is a genuine
marginal information only when $\Imat_{MM}$ is full-rank; for a rank-deficient $\Imat_{MM}$
the pseudo-inverse makes $\bm S$ a \emph{regularized surrogate} that need not be PSD (a
negative marginal eigenvalue in an early single-beacon fixture was exactly this artifact).
All networks reported here have full-rank, positive-definite $M$-blocks, so $\bm S$ is a true
marginal and no reported value is affected. The eigenvalues of $\bm S$ (via
\code{fim::sym\_eig}, ascending) are reported as \code{marginal\_eigs}; $\lambda_{\min}(\bm S)$
is the coupling-1 near-defect strength. For the \code{well\_posed} network,
$\bm S$ is positive-definite with eigenvalues $[1.8720,\,35.9995]$.

\subsection{Independence of the numpy/LAPACK oracle}
\label{app:oracle}

The oracle \code{scripts/gen\_coupled\_gauge\_ref.py} reproduces the pipeline through a
\emph{different} numerical stack: eigenvalues by \code{np.linalg.eigvalsh}, null-space Gram
ranks by \code{np.linalg.eigh}, and the Schur pseudo-inverse by \code{np.linalg.pinv} (SVD),
against the engine's Jacobi eigensolver (\code{fim::sym\_eig}) and spectral-sum pseudo-inverse
(\code{fim::crlb}). The two backends---Jacobi versus LAPACK divide-and-conquer, spectral-sum
versus SVD pseudo-inverse---are independent implementations, so agreement to relative and
absolute error below $10^{-3}$ is a genuine cross-check, not a tautology. Non-circularity is
structural: the oracle reads the observation \emph{rows} from \code{network.json} and derives
\code{reference.json} from them, never from a Rust output; six-decimal rounding of the rows
before both writing and reading guarantees byte-identical inputs.

\subsection{Derivation of the frame\texorpdfstring{$\to$}{->}rate Jacobian}
\label{app:jacobian}

The three coefficients \eqref{eq:dads}--\eqref{eq:dadr} follow from the post-Newtonian clock
rate. We adopt the convention---used once here and throughout---that $\alpha$ is the rate of
a lunar-surface clock measured \emph{relative to Terrestrial Time} (TT), with TT tied to the
defining Earth geopotential $W_{0,\oplus}$. To first order a clock's rate is set by the
gravitational redshift and the second-order Doppler (kinetic) term, so
\begin{equation}
  \alpha \;\simeq\; \frac{W_{0,\oplus}-\Umoon}{c^2}-\frac{v^2}{2c^2}+\text{const},
  \label{eq:rate}
\end{equation}
with $\Umoon>0$ the lunar self-potential and $v$ the speed
\citep{soffel2003iau,mueller2018clocks}. In this surface-relative-to-TT convention the
self-potential enters with a \emph{minus} sign---a deeper lunar well runs the clock slower
relative to TT---so the coefficients below carry the $+\Umoon/c^2$ and $+\gmoon/c^2$ signs of
the $(W_0-U)/c^2$ form, opposite the sign of the bare redshift $\Phi=-GM/r$. For a clock on
the lunar surface the self-potential is $\Umoon=GM_{\mathrm{M}}/\rmoon$ and the surface
gravity is $\gmoon=GM_{\mathrm{M}}/\rmoon^{2}$. Differentiating:
\begin{itemize}
  \item \textbf{Scale.} A fractional scale change $\delta s$ scales the lunar radius,
  $\rmoon\mapsto\rmoon(1+\delta s)$, so the surface self-potential shrinks,
  $\Umoon(s)=\Umoon(1-s)$ to first order. Substituting into
  $\alpha=(W_{0,\oplus}-\Umoon(s))/c^2-v^2/(2c^2)+\text{const}$ gives
  $\partial\alpha/\partial s=+\Umoon/c^2=+3.140\times10^{-11}$ per unit fractional scale, with
  $GM_{\mathrm{M}}=4.9028\times10^{12}\,\mathrm{m^3/s^2}$ and
  $\rmoon=1.7374\times10^{6}$~m.
  \item \textbf{Radial.} A radial displacement $\delta r$ changes the potential by
  $\gmoon\,\delta r$, so $\partial\alpha/\partial r=+\gmoon/c^2=+1.807\times10^{-17}$ per
  metre. The identity $\partial\alpha/\partial r=(\partial\alpha/\partial s)/\rmoon$ is exact
  ($\gmoon/c^2=(\Umoon/\rmoon)/c^2$) and unit-tested to $10^{-12}$.
  \item \textbf{Velocity.} The kinetic term gives $\partial\alpha/\partial v=-|v|/c^2$, with
  $|v|$ the geocentric Moon speed at the epoch (from \code{lunar\_time}); at 2024-01-01
  ($|v|\approx969$~m/s) this is $-1.078\times10^{-14}$, near J2000 ($|v|\approx997$~m/s)
  $\approx-1.11\times10^{-14}$ per (m/s). The velocity coefficient is epoch-dependent by the
  several-percent monthly variation of the Moon's speed (the quoted $969$ vs $997$~m/s
  already differ by nearly $3\%$).
\end{itemize}
The external (Earth) tidal potential is \emph{excluded} from $\Umoon$: by the equivalence
principle a freely-falling body does not sense the uniform part of the external potential, so
$GM_{\oplus}/d$ does not enter the clock rate; only the tidal \emph{gradient} does, at
$\sim2.7\times10^{-22}$ per metre, negligible. All three coefficients are \Modelled{}
(first-principles, cross-checked to the $[56,59]\,\mu$s/day lunar-clock-rate figures
% confirm the [56,59] $\mu$s/day figure against fienga2024moonlight at camera-ready
consistent with the Moonlight reference-system literature \citep{fienga2024moonlight}), not
certified operational values.
